%BeginMC
\question
During protist conjugation, which of the following does \emph{not} occur during some part of the process?

\begin{choices}
\correctchoice The macronucleus fuses with the micronucleus. %EndChoice
\choice The micronucleus undergoes meiosis. %EndChoice
\choice Cells exchange micronuclei.  %EndChoice
\choice The macronucleus distintegrates. %EndChoice
\choice The micronucleus undergoes mitosis. %EndChoice
\choice Some micronuclei distintegrate. %EndChoice
\end{choices}
%EndMC


%BeginMC
\question
Hydrostatic skeletons are 

\begin{choices}
\correctchoice fluid-filled compartments under pressure. %EndChoice
\choice external skeletons that provide support and protection against dehydration. %EndChoice
\choice internal skeletons that are buried in soft tissue.  %EndChoice
\choice skeletons found in deep sea organisms that depend on high pressure of the surrounding water to provide support. %EndChoice
\choice skeletons made of microtubules in the cytoplasm of protists. %EndChoice
\end{choices}
%EndMC

%BeginMC
\question
Which of the following are primarily responsible for movement in the eukaryotic flagellum?

\begin{choices}
\correctchoice Dynein and microtubules. %EndChoice
\choice Basal apparatus and hook. %EndChoice
\choice Actin and myosin filaments.  %EndChoice
\choice Hydrogen ions and basal apparatus. %EndChoice
\choice Sensory and motor neurons. %EndChoice
\end{choices}
%EndMC

%BeginMC
\question
Differences in energy quality and availability affects the amount of surface area and volume in organisms. Compared to animals, plants have 

\begin{choices}
\correctchoice high surface area and low volumes. %EndChoice
\choice high surface area and high volumes. %EndChoice
\choice low surface areas and low volumes.  %EndChoice
\choice low surface area and high volumes. %EndChoice
\end{choices}
%EndMC

%BeginMC
\question
Ignoring all other factors, what kind of day would result in the fastest delivery of water and minerals to the leaves of a tree?

\begin{choices}
\correctchoice A warm, dry day. %EndChoice
\choice A cool, dry day. %EndChoice
\choice A warm, humid day.  %EndChoice
\choice A cool, humid day. %EndChoice
\choice A very hot, dry, windy day. %EndChoice
\end{choices}
%EndMC

%BeginMC
\question
Most of the water taken up by a plant

\begin{choices}
\correctchoice is lost during transpiration. %EndChoice
\choice is converted to sugars during photosynthesis. %EndChoice
\choice enters the central vacuoles of cells.  %EndChoice
\choice keep guard cells turgid. %EndChoice
\choice transports sugars in the phloem. %EndChoice
\end{choices}
%EndMC

%BeginMC
\question
Transpiration in plants requires all of the following \emph{except}

\begin{choices}
\correctchoice active transport through xylem cells. %EndChoice
\choice adhesion of water molecules to cellulose. %EndChoice
\choice cohesion between water molecules.  %EndChoice
\choice evaporation of water molecules. %EndChoice
\choice water vapor on the spongy mesophyll. %EndChoice
\end{choices}
%EndMC


%BeginMC
\question
Galvanotaxis occurs when 

\begin{choices}
\correctchoice protists respond to an electrical field. %EndChoice
\choice animals respond to metal ions. %EndChoice
\choice plants respond to gravity.  %EndChoice
\choice animals respond to temperature. %EndChoice
\choice protists respond to touch. %EndChoice
\choice plants respond to hormones. %EndChoice
\end{choices}
%EndMC


%BeginMC
\question
In terms of structure, cilia  are most like 
\begin{choices}
\correctchoice eukaryotic flagella. %EndChoice
\choice bacterial flagella. %EndChoice
\choice microtubules.  %EndChoice
\choice actin filaments. %EndChoice
\choice myosin filaments. %EndChoice
\end{choices}
%EndMC

%BeginMC
\question
When ATP attaches to the myosin head

\begin{choices}
\correctchoice the myosin head releases from the actin filament. %EndChoice
\choice the actin head attaches to the myosin filament. %EndChoice
\choice the myosin head attaches to the actin filament.  %EndChoice
\choice the actin head attaches to the myosin filament. %EndChoice
\choice the myosin head initiates the power stroke. %EndChoice
\end{choices}
%EndMC

%BeginMC
\question
The ion used most often as a second messenger is

\begin{choices}
\correctchoice \ce{Ca+} %EndChoice
\choice \ce{Na+} %EndChoice
\choice \ce{K+} %EndChoice
\choice \ce{H+} %EndChoice
\choice \ce{Cl-} %EndChoice
\end{choices}
%EndMC

%BeginMC
\question
The main stimulus used by birds to find their way north for the summer time is 
\begin{choices}
\correctchoice electomagnetic. %EndChoice
\choice thermal. %EndChoice
\choice chemical. %EndChoice
\choice mechanical. %EndChoice
\choice thigmic. %EndChoice
\end{choices}
%EndMC


%BeginMC
\question
The vascular tissue system for delivery of water, minerals, and nutrients throughout the plant includes all of the following except for

\begin{choices}
\correctchoice cohesion. %EndChoice
\choice sieve plates. %EndChoice
\choice phloem. %EndChoice
\choice xylem. %EndChoice
\choice companion cells. %EndChoice
\end{choices}
%EndMC


%BeginMC
\question
Pores on the leaf surface that function in gas and water exchange are called

\begin{choices}
\correctchoice stomata. %EndChoice
\choice waxy cuticles. %EndChoice
\choice spongy mesophyll. %EndChoice
\choice xylem. %EndChoice
\choice phloem. %EndChoice
\end{choices}
%EndMC


%BeginMC
\question
What drives the flow of water through the xylem?

\begin{choices}
\correctchoice The evaporation of water from the leaves. %EndChoice
\choice Transpirational pull through sieve-tube elements. %EndChoice
\choice The number of companion cells in the phloem.  %EndChoice
\choice Abscisic acid opening guard cells around the stomata. %EndChoice
\choice A higher concentration of sugars in the leaves than in the roots. %EndChoice
\end{choices}
%EndMC


%BeginMC
\question
From an energy perspective, animals are mobile because they consume

\begin{choices}
\correctchoice low quality energy and carbon intake is irregular.%EndChoice
\choice high quality energy and carbon intake is constant. %EndChoice
\choice low quality energy and low energy efficiency.%EndChoice
\choice high quality energy and high energy efficiency. %EndChoice
\choice high quality energy and carbon intake is irregular. %EndChoice
\choice low quality energy and carbon intake is constant. %EndChoice
\end{choices}
%EndMC


%BeginMC
\question
Which of the following about phloem is false?

\begin{choices}
\correctchoice The plasma membrane is modified to form a sieve plate. %EndChoice
\choice Sieve-tube elements do not have a nucleus. %EndChoice
\choice Sugars move by passive diffusion. %EndChoice
\choice The cell has cytoplasm. %EndChoice
\choice Companion cells provide proteins to sieve-tube elements. %EndChoice
\end{choices}
%EndMC


%BeginMC
\question
The highest quality form of energy (best able to do work) for organisms is

\begin{choices}
\correctchoice radiant. %EndChoice
\choice thermal. %EndChoice
\choice chemical. %EndChoice
\choice kinetic. %EndChoice
\choice potential. %EndChoice
\end{choices}
%EndMC

%BeginMC
\question
The protein that is activated in response to auxin that allows phototropism to occur is

\begin{choices}
\correctchoice expansin. %EndChoice
\choice proton channel. %EndChoice
\choice transmembrane. %EndChoice
\choice abscisic acid. %EndChoice
\choice wallulase. %EndChoice
\choice basal apparatus. %EndChoice
\end{choices}
%EndMC

%BeginMC
\question
The fluid that moves around in the circulatory system of a typical arthropod is the

\begin{choices}
\correctchoice hemolymph. %EndChoice
\choice gastrovascular fluid. %EndChoice
\choice blood plasma. %EndChoice
\choice cytoplasm. %EndChoice
\choice intracellular fluid. %EndChoice
\end{choices}
%EndMC

%BeginMC
\question
Organisms with a circulating fluid that is separate from the fluid that directly surrounding the body's cells are likely to have

\begin{choices}
\correctchoice a closed circulatory system. %EndChoice
\choice an open circulatory system. %EndChoice
\choice a gastrovascular cavity. %EndChoice
\choice a hemolymph system. %EndChoice
\choice a vascular system. %EndChoice
\end{choices}
%EndMC

%BeginMC
\question
The function of valves in veins in a closed circulatory system is to 
\begin{choices}
\correctchoice ensure one way flow of circulatory fluids. %EndChoice
\choice decrease the amount of pressure in the veins.  %EndChoice
\choice provide support for the extra muscles around the vein. %EndChoice
\choice enable the vein to withstand high pressure. %EndChoice
\choice enhance diffusion of nutrients to the tissues. %EndChoice
\end{choices}
%EndMC

%BeginMC
\question
The type of nerve cell responsible for integrating different stimuli is the 

\begin{choices}
\correctchoice interneuron %EndChoice
\choice dendrite %EndChoice
\choice motor neuron %EndChoice
\choice sensory neuron %EndChoice
\choice axon %EndChoice
\end{choices}
%EndMC

